% vim: ai sw=2 spell spelllang=cs

\begin{frame}
  \frametitle{Komunikace mezi procesy}

  \heading{LDD3 část kap. 6 (zastaralá)}

  \iffimuni
    \bigskip

    \tableofcontents
  \fi

\end{frame}

\mysection{Čekání na událost}

\iffimuni
\subsection{Completion}
\fi

\begin{frame}[fragile]
  \frametitle{Jednoduchá komunikace}

  \begin{itemize}
    \item 2 procesy (producent-konzument)
    \begin{itemize}
      \item A: čeká na nějakou hotovou práci
      \item B: udělá nějakou práci a oznámí A dokončení
      \item A: pokračuje
    \end{itemize}
  \end{itemize}

  \pause

  \heading{API}
  \begin{itemize}
    \item \header{linux/completion.h}, \code{struct completion}
    \item \code{DECLARE_COMPLETION}$_S$, \code{init_completion}$_D$
    \item Čekání: \code{wait_for_completion},
      \code{wait_for_completion_interruptible} (retval)
    \item Dokončení: \code{complete}, \code{complete_all}
  \end{itemize}

  \pause

  \begin{block}{Příklad}
  \begin{lstlisting}
static DECLARE_COMPLETION(my_comp);
static char str[32];
A                                        B
wait_for_completion(&my_comp);           strcpy(str, "Ahoj");
pr_info("%s\n", str);                    complete(&my_comp);
  \end{lstlisting} %stop_zone
  \end{block}

\end{frame}

\begin{frame}
  \frametitle{Úkol}

  \heading{Použití \code{completion}}
  \begin{enumerate}
    \item Do \file{pb173/11/events} doplňte \code{completion}
    \item Číst se bude, až někdo dokončí zápis
    \begin{itemize}
      \item Čekání v \code{read} bude možné přerušit signálem
	(\code{wait_for_completion_interruptible})
      \item V případě signálu vraťte \code{\-ERESTARTSYS}
    \end{itemize}

    \item Spusťte 2 instance \cmd{cat /dev/my_name}
    \item Spusťte několikrát \cmd{echo XXX > /dev/my_name}

    \item Pozorujte, co se děje, když použijete
    \begin{itemize}
      \item \code{complete} a potom:
      \item \code{complete_all}
    \end{itemize}
  \end{enumerate}

\end{frame}

\iffimuni
\subsection{Wait event a ostatní}
\fi

\begin{frame}
  \frametitle{Složitější komunikace}

  \begin{itemize}
    \item Čekání na více místech na jinou událost
    \item 3 procesy
    \begin{itemize}
      \item A: plní buffer
      \item B: čeká na 100\,B
      \item C: čeká na 200\,B
    \end{itemize}
  \end{itemize}

  \pause

  \heading{API}
  \begin{itemize}
    \item \header{linux/wait.h}, \code{wait_queue_head_t}
    \item \code{DECLARE_WAIT_QUEUE_HEAD}$_S$,
      \code{init_waitqueue_head}$_D$
    \item Čekání: \code{wait_event}, \code{wait_event_interruptible}
    \item Dokončení: \code{wake_up}, \code{wake_up_all}
  \end{itemize}

  \pause
  \bigskip

  Pozn.: \code{completion} je jen \code{wait_queue} + počítadlo

\end{frame}

\begin{frame}[fragile]{Složitější komunikace}{Příklad}

  \begin{lstlisting}
static DECLARE_WAIT_QUEUE_HEAD(my_wait);
static atomic_t my_cnt = ATOMIC_INIT(0);
static char str[512];
...
A
  for (a = 0; a < sizeof(str); a++) {
    str[atomic_inc_return(&my_cnt) - 1] = 'A';
    wake_up_all(&my_wait);
    msleep(50);
  }
  \end{lstlisting}

  \pause

  \begin{lstlisting}
B
  wait_event(my_wait, atomic_read(&my_cnt) > 100);
  pr_info("%s\n", str);
  \end{lstlisting} %stop_zone

  \pause

  \begin{lstlisting}
C
  if (wait_event_interruptible(my_wait, atomic_read(&my_cnt) > 200)) {
    pr_info("interrupted\n");
    return;
  }
  pr_info("%s\n", str);
  \end{lstlisting} %stop_zone

\end{frame}

\begin{frame}
  \frametitle{Úkol}

  \heading{Použití \code{wait_queue}}
  \begin{enumerate}
    \item V \file{pb173/11/events} použijte namísto \code{completion}
      \code{wait_queue}
    \item Číst se bude, až bude v bufferu zapsáno alespoň 5 znaků
  \end{enumerate}

\end{frame}

\iffimuni
\subsection{Scheduler a probouzení procesů}
\fi

\begin{frame}
  \frametitle{Scheduler}

  \begin{itemize}
    \item Celé je to instruování plánovače (\header{linux/sched.h})

    \item Ve skutečnosti je \code{wait_event}:
    \begin{itemize}
      \item Nastavení typu spícího stavu procesu: \code{set_current_state}
      \begin{itemize}
	\item D stav: \code{TASK_UNINTERRUPTIBLE}
	\item S stav: \code{TASK_INTERRUPTIBLE}
      \end{itemize}

      \pause

      \item Uspání procesu bez/s timeoutem:
        \code{schedule}/\code{schedule_timeout}
    \end{itemize}

    \pause
    
    \item A \code{wake_up}:
    \begin{itemize}
      \item Probuzení procesu: \code{wake_up_process}
      \item Popř.\ s obsluhou signálů: \code{signal_pending}
    \end{itemize}

    \pause

    \item \code{completion} a \code{wait_queue} ulehčuje práci
    \begin{itemize}
      \item Pamatováním si, koho vzbudit
      \item Případnou kontrolou signálů
      \item Případnou kontrolou vypršení timeoutu
      \item \code{completion} má navíc čítač, kolik procesů vzbudit
    \end{itemize}
  \end{itemize}

\end{frame}

\begin{frame}
  \frametitle{Úkol}

  \heading{Explicitní čekání}
  \begin{enumerate}
    \item Místo \code{wait_event}, použijte funkce z předchozího slidu
    \item Jako vzor prostudujte kód ve funkci \code{afu_read}

    \item Zkopírujte jej a použijte
    \begin{itemize}
      \item Odstraňte zámek a ostatní irelevantní kód
    \end{itemize}

    \item Pozměňte čekací podmínku
    \begin{itemize}
      \item Bude se čekat alespoň na 2~,,\code{a}'' a 1~,,\code{b}'' v poli
    \end{itemize}
  \end{enumerate}

\end{frame}

\mysection{Vlákna}

\begin{frame}[fragile]
  \frametitle{\code{kthread}}

  \begin{itemize}
    \item Vytvoření vlákna (procesu) pro výpočty
    \begin{itemize}
      \item Běží v jádře
      \item Chování jako uživatelský proces
    \end{itemize}

    \pause

    \item \header{linux/kthread.h}
    \item Vytvoření, zrušení: \code{kthread_run}, \code{kthread_stop}
    \item Test na skončení ve vlákně: \code{kthread_should_stop}
  \end{itemize}

  \pause

  \begin{block}{Příklad}
  \begin{center}
  \begin{tabular}{c|c}
  \begin{lstlisting}
static int my_fun(void *data)
{ /* data = my_data */
  while (1) {
    wait_event_interruptible(...);
    if (kthread_should_stop())
      break;
    do_some_work();
  }
  return 0;
}
  \end{lstlisting}~ &
  \begin{uncoverenv}<4->
  ~\begin{lstlisting}
struct task_struct *my_t;
static int my_init(void)
{
  my_t = kthread_run(my_fun, my_data, "my_fun");
  return IS_ERR(my_t) ? PTR_ERR(my_t) : 0;
}
static void my_exit(void)
{
  kthread_stop(my_t);
}
  \end{lstlisting}
  \end{uncoverenv}
  \end{tabular}
  \end{center}
  \end{block}

\end{frame}

\begin{frame}
  \frametitle{Úkol}

  \heading{Vytvoření vlákna}
  \begin{enumerate}
    \item Vytvořte vlákno
    \item Vlákno počká na \code{write} až někdo něco zapíše
    \item Přepíše všechny znaky ,,\code{a}'' na ,,\code{b}''
    \item Vzbudí všechny v \code{read}, kteří čekají na data
    \item Opakujte v cyklu dokud někdo neodebere modulu
    \item Vyzkoušejte zápisem a čtením do/z \file{/dev/}
  \end{enumerate}

\end{frame}
